\appendix
\section{Claim-level evidence ledger}
Status and design code are separate. An E0 citation supplies background or provenance without establishing the whole bridge.
\subsection*{C01 · Specific constitutional variants \ensuremath{\rightarrow} tissue phenotype}
\noindent\textbf{Supported · E1 · Human CAH-X genotype–phenotype cohort}\par
\textit{Claim.} CYP21A2–TNXB disruption supplies a defined endocrine/connective-tissue interface.\par
\textit{Limit.} This does not validate the broad RCCX neurodevelopmental cluster. The autonomic frequency in selected hEDS/HSD clinics cannot be generalized to all carriers.\par
\textit{Rival.} Other genetic and clinical pathways explain similar tissue findings outside CAH-X.\par
\textit{Discriminator.} A prespecified RCCX contribution must add replicated prediction beyond broader constitutional measures.\par
\textit{Sources.} \cite{ref13,projP1,ref22,ref33}\par
\subsection*{C02 · Birth conditions \ensuremath{\rightarrow} AVP / copeptin response}
\noindent\textbf{Supported · E2 · Human newborn cord blood}\par
\textit{Claim.} Delivery-associated AVP/copeptin differences are reported across human studies.\par
\textit{Limit.} An acute peripheral signal does not establish persistent neural calibration; OXT is not interchangeable.\par
\textit{Rival.} Obstetric indication and the acute transition explain the cord signal without lasting calibration.\par
\textit{Discriminator.} A persistent-developmental claim needs longitudinal effects beyond the transient birth response.\par
\textit{Sources.} \cite{ref1,ref2}\par
\subsection*{C03 · Maternal ecology \ensuremath{\rightarrow} infant strain inheritance}
\noindent\textbf{Supported · E1 · Human mother–infant longitudinal cohort}\par
\textit{Claim.} Maternal gut reservoirs, delivery, feeding and subsequent contacts shape measured strain acquisition.\par
\textit{Limit.} Delivery label is a proxy; observational transmission does not establish a behavioral effect.\par
\textit{Rival.} Delivery mode is a marker for feeding, maternal reservoirs, antibiotics, and household ecology.\par
\textit{Discriminator.} The nominated strain/function fails to add held-out prediction of prespecified outcomes.\par
\textit{Sources.} \cite{ref3}\par
\subsection*{C04 · \emph{L. reuteri} \ensuremath{\rightarrow} vagus / OXT / VTA social plasticity}
\noindent\textbf{Established · E3 · Mouse ASD-related models}\par
\textit{Claim.} Vagal integrity and oxytocin signaling participate in experimentally observed social-behavior rescue.\par
\textit{Limit.} The complete pathway has not been causally demonstrated in humans.\par
\textit{Rival.} Other microbial metabolites or immune routes mediate the mouse social effect.\par
\textit{Discriminator.} Replicated pathway-specific disruption fails to alter the effect in the nominated preparation.\par
\textit{Sources.} \cite{ref4}\par
\subsection*{C05 · \emph{L. reuteri} formulation \ensuremath{\rightarrow} selected human outcomes}
\noindent\textbf{Supported · E4 · Pilot trials: children; adolescents/adults}\par
\textit{Claim.} Human trials provide limited social or adaptive-behavior signals.\par
\textit{Limit.} Global autism severity is not consistently improved; human OXT/vagus mediation remains unproven.\par
\textit{Rival.} Trial-specific formulation effects, chance, or multiple endpoints explain the selected benefit.\par
\textit{Discriminator.} Adequately powered replication fails on prespecified social outcomes.\par
\textit{Sources.} \cite{ref5,ref6}\par
\subsection*{C06 · Thrombospondin–\ensuremath{\alpha_2\delta}-1 \ensuremath{\rightarrow} synaptogenesis}
\noindent\textbf{Established · E3 · Rodent neural preparations and developing mouse cortex}\par
\textit{Claim.} \ensuremath{\alpha_2\delta}-1 participates in excitatory synapse formation; postsynaptic Rac1 is mechanistically implicated.\par
\textit{Limit.} Not identical to presynaptic trafficking; not a demonstrated consequence of birth mode. A later binding assay found soluble TSP4–\ensuremath{\alpha_2\delta}-1 binding was gabapentin-insensitive; binding and synapse outcomes must be separated.\par
\textit{Rival.} Thrombospondin signaling acts through downstream partners independently of CaV trafficking.\par
\textit{Discriminator.} Orthogonal replication fails for the specified receptor-dependent synaptogenic effect.\par
\textit{Sources.} \cite{ref10,ref11,ref31}\par
\subsection*{C07 · \ensuremath{\alpha_2\delta} maturation \ensuremath{\rightarrow} release regulation}
\noindent\textbf{Established · E3 · Experimentally manipulated neuronal culture}\par
\textit{Claim.} Proteolytic maturation of \ensuremath{\alpha_2\delta} affects channel association and vesicular release probability.\par
\textit{Limit.} No verified human perinatal cleavage abnormality is inferred.\par
\textit{Rival.} Other maturation proteins or activity-dependent changes explain release regulation.\par
\textit{Discriminator.} Selective maturation manipulation fails to reproduce the release effect in the claimed preparation.\par
\textit{Sources.} \cite{ref12}\par
\subsection*{C08 · Development \ensuremath{\rightarrow} channel contribution reallocation}
\noindent\textbf{Established · E3 · Rat auditory brainstem, postnatal days 4–14}\par
\textit{Claim.} Presynaptic channel contributions change during development in the studied synapse.\par
\textit{Limit.} Not a universal brain switch, a human-age calibration, or proof of arrested development. Mouse expression profiling found stage- and region-specific programs; one activity-blockade condition did not alter the measured expression profile.\par
\textit{Rival.} Intrinsic developmental programs explain the channel shift without the nominated exposure.\par
\textit{Discriminator.} Within-preparation replication contradicts the specified temporal change.\par
\textit{Sources.} \cite{ref8,ref30}\par
\subsection*{C09 · Perinatal signals \ensuremath{\rightarrow} \ensuremath{\alpha_2\delta}–CaV2 state}
\noindent\textbf{Proposed · E0 · Target developmental stage must be specified}\par
\textit{Claim.} Endocrine, immune or microbial signals might alter channel organization and synaptogenesis.\par
\textit{Limit.} This is the main missing bridge. Molecular plausibility does not demonstrate it.\par
\textit{Rival.} Early exposure affects networks through a pathway that bypasses \ensuremath{\alpha_2\delta}/CaV2.\par
\textit{Discriminator.} Precise null exposure effects, or failed mediator disruption/rescue, weaken the nominated bridge.\par
\textit{Sources.} \cite{projP3,ref10,ref12}\par
\subsection*{C10 · Network calibration \ensuremath{\rightarrow} regulatory phenotype}
\noindent\textbf{Proposed · E0 · Cross-level developmental hypothesis}\par
\textit{Claim.} Synaptic organization may influence later sensory, autonomic and affiliative regulation.\par
\textit{Limit.} The integrated mapping lacks direct longitudinal validation; many bypasses are possible.\par
\textit{Rival.} Ongoing sleep, pain, reward and social learning explain later regulation without an early molecular mediator.\par
\textit{Discriminator.} The specified network measures fail to predict or mediate the targeted outcome.\par
\textit{Sources.} \cite{projP1,projP3,ref9,ref14,ref15}\par
\subsection*{C11 · Regulatory phenotype \ensuremath{\times} environment \ensuremath{\rightarrow} closure}
\noindent\textbf{Horizon · E0 · Proposed repeated human task measures}\par
\textit{Claim.} RCB compares positive external and internal regulatory gains on the same scale.\par
\textit{Limit.} Not a validated trait, diagnosis, or normative scoring tool.\par
\textit{Rival.} Conventional arousal, executive function or reward sensitivity explains task-dependent gains.\par
\textit{Discriminator.} Poor reliability or no incremental out-of-sample interaction beyond conventional measures.\par
\textit{Sources.} \cite{projP1,projP2,ref16}\par
\subsection*{C12 · Regulation + context \ensuremath{\rightarrow} developmental reorganization}
\noindent\textbf{Horizon · E0 · Dąbrowski-inspired developmental interpretation}\par
\textit{Claim.} Different support and load trajectories may produce recovery, reorganization or persistent impairment.\par
\textit{Limit.} The bifurcation is not an established biological stage model; outcomes are not moral ranks.\par
\textit{Rival.} Resources, support and treatment explain trajectories without a discrete reorganization process.\par
\textit{Discriminator.} Operationalized trajectories fail to replicate or the construct adds no value over alternatives.\par
\textit{Sources.} \cite{projP1,ref17}\par
\subsection*{C13 · Human–AI coupling \ensuremath{\rightarrow} conversion and dissemination}
\noindent\textbf{Horizon · E0 · Proposed task- and resource-dependent interaction}\par
\textit{Claim.} Responsive tools may reduce implementation costs but can increase branching and dependence.\par
\textit{Limit.} AI writing evidence does not establish an RCB-specific advantage or civilizational inversion.\par
\textit{Rival.} Generic AI assistance and unequal opportunity explain output changes without phenotype-specific selection.\par
\textit{Discriminator.} Independent RCB fails to predict the nominated environment interaction on verified completed output.\par
\textit{Sources.} \cite{projP2,ref21}\par
\subsection*{C14 · Early immune ecology \ensuremath{\rightarrow} disgust sensitivity calibration}
\noindent\textbf{Horizon · E0 · Unverified developmental bridge}\par
\textit{Claim.} Early immune/microbial conditions might influence later pathogen avoidance.\par
\textit{Limit.} Adult disgust associations do not establish early microbial causation or political consequences.\par
\textit{Rival.} Current learning and culture explain disgust variation without early microbial calibration.\par
\textit{Discriminator.} No replicated prospective relationship with age-appropriate germ-aversion measures.\par
\textit{Sources.} \cite{projP3,ref7}\par
\subsection*{C15 · Phenotype \ensuremath{\rightarrow} selected environment \ensuremath{\rightarrow} later phenotype}
\noindent\textbf{Proposed · E0 · Time-indexed feedback hypothesis}\par
\textit{Claim.} Actions change subsequent exposure, reinforcement and support, feeding into later regulation.\par
\textit{Limit.} A cyclic diagram is not an identified causal model; time-order and confounding matter.\par
\textit{Rival.} Stable traits and environments jointly cause behavior and exposure without strong selection feedback.\par
\textit{Discriminator.} Time-indexed interventions and predictions fail against a nonfeedback alternative.\par
\textit{Sources.} \cite{projP1,projP2}\par
\subsection*{C16 · Birth mode \ensuremath{\rightarrow} categorical ASD / ADHD}
\noindent\textbf{Supported · E1 · Human observational counterevidence}\par
\textit{Claim.} Fully adjusted NHS II estimates are compatible with no association.\par
\textit{Limit.} Null categorical effects constrain broad claims; they do not test every dimensional mediator.\par
\textit{Rival.} Confounding explains crude birth-mode associations; adjusted categorical effects are near null.\par
\textit{Discriminator.} A well-identified replicated causal estimate would change this assessment.\par
\textit{Sources.} \cite{ref19}\par
\subsection*{C17 · VMT assignment \ensuremath{\rightarrow} 12-month social-emotional score}
\noindent\textbf{Supported · E4 · Pilot infant trial}\par
\textit{Claim.} Randomized VMT versus saline favored a social-emotional score; global ASQ-3 was not significant.\par
\textit{Limit.} No identified strain or \ensuremath{\alpha_2\delta} mediator; vaginal reference group was not randomized.\par
\textit{Rival.} The intervention affects social scores through a pathway independent of the named organism or molecular hinge.\par
\textit{Discriminator.} Larger preregistered replication fails to confirm the specific outcome.\par
\textit{Sources.} \cite{ref20}\par
\subsection*{C18 · RCCX susceptibility \ensuremath{\rightarrow} \ensuremath{\alpha_2\delta} hinge}
\noindent\textbf{Proposed · E0 · No direct evidence located in this pass}\par
\textit{Claim.} Constitution could modify sensitivity to developmental perturbation.\par
\textit{Limit.} Specific RCCX-to-channel mediation is unproven; genomic proximity is insufficient. ADHD GWAS is polygenic; a C4 schizophrenia result does not establish this neurodevelopmental bridge.\par
\textit{Rival.} Distributed genetic architecture and independent perinatal factors explain the data without an RCCX–hinge link.\par
\textit{Discriminator.} No prespecified RCCX interaction or incremental prediction in independent cohorts.\par
\textit{Sources.} \cite{projP1,projP3,ref13,ref22,ref23,ref24}\par
\subsection*{C19 · Phenibut pharmacology \ensuremath{\rightarrow} hypothesis generation}
\noindent\textbf{Established · E3 · Experimental binding and animal pharmacology}\par
\textit{Claim.} R-phenibut has \ensuremath{\alpha_2\delta} and GABA-B pharmacological activity.\par
\textit{Limit.} Adult response does not diagnose a target-specific or developmental lesion.\par
\textit{Rival.} GABA-B effects, expectancy, or current state explain an adult response without developmental \ensuremath{\alpha_2\delta} pathology.\par
\textit{Discriminator.} The developmental inference is not identified by a response; it requires independent tests.\par
\textit{Sources.} \cite{ref18,projP3}\par
\subsection*{C20 · Constitution \ensuremath{\times} perinatal conditions \ensuremath{\rightarrow} initialization}
\noindent\textbf{Proposed · E0 · Unverified integrated developmental interaction}\par
\textit{Claim.} Measured constitutional factors may change responses to perinatal exposures.\par
\textit{Limit.} The broad interaction is not established by CAH-X or molecular synapse studies.\par
\textit{Rival.} Constitution and perinatal conditions contribute additively, with no specific interaction.\par
\textit{Discriminator.} Prespecified constitutional effect modification fails in adequately precise independent studies.\par
\textit{Sources.} \cite{projP1,projP3,ref13}\par
\subsection*{C21 · Early adversity \ensuremath{\leftrightarrow} stress-regulatory methylation}
\noindent\textbf{Supported · E1 · Human postmortem and blood observational studies; tissue and age differ}\par
\textit{Claim.} Adversity history is associated with some NR3C1 or FKBP5 methylation and expression measures in selected samples.\par
\textit{Limit.} A 113-dyad study found no direct transfer at assayed newborn sites, and a 3,965-person population study did not replicate a maltreatment interaction. This cannot identify a universal inherited mark.\par
\textit{Rival.} Genotype, cell mixture, current illness, shared environment, and selection explain the associations.\par
\textit{Discriminator.} Prospective, cell-composition-aware replication should show the nominated site and direction across independent cohorts.\par
\textit{Sources.} \cite{ref15,ref25,ref26,ref34}\par
\subsection*{C22 · Task event rate \ensuremath{\rightarrow} attention performance}
\noindent\textbf{Supported · E2 · Go/No-Go ADHD task meta-analysis; not a general personality scale}\par
\textit{Claim.} Slow event rates disproportionately slowed ADHD groups; fast rates increased commission-error effects in the reviewed tasks.\par
\textit{Limit.} General ADHD–control differences persisted and reaction-time variability did not show the predicted event-rate effect.\par
\textit{Rival.} Specific task learning or reward expectations, rather than regulatory closure, explain the pattern.\par
\textit{Discriminator.} Matched within-person task crossovers fail to reproduce event-rate interactions beyond task demands.\par
\textit{Sources.} \cite{ref27}\par
\subsection*{C23 · Camouflaging is measurable; broad masking remains a bridge}
\noindent\textbf{Supported · E1 · Autistic adult psychometric samples}\par
\textit{Claim.} The CAT-Q measured reported social camouflaging in autistic and non-autistic adults.\par
\textit{Limit.} This does not prove unconscious homeostatic masking across ADHD, trauma, or other phenotypes.\par
\textit{Rival.} Conventional social adaptation and context account for reported strategies.\par
\textit{Discriminator.} Longitudinal behavioral and physiological measures fail to distinguish generalized masking from social learning.\par
\textit{Sources.} \cite{ref32}\par
\subsection*{C24 · Chronotype diversity \ensuremath{\rightarrow} group-level nighttime coverage}
\noindent\textbf{Supported · E1 · Hadza actigraphy, 20 days; no ADHD measure}\par
\textit{Claim.} At least one group member was scored awake or in light sleep in 99.8\% of sampled nighttime epochs.\par
\textit{Limit.} The result supports group-level heterogeneity, not an ADHD subtype or general adaptive superiority.\par
\textit{Rival.} Age structure and ordinary awakenings explain the coverage without a specialized vigilance trait.\par
\textit{Discriminator.} Comparable groups with mixed age and chronotype fail to show the predicted coverage.\par
\textit{Sources.} \cite{ref28}\par
\subsection*{C25 · Hypermobility cohorts show heterogeneous autonomic profiles}
\noindent\textbf{Supported · E1 · Selected 102-person hEDS/HSD adult clinical cohort}\par
\textit{Claim.} Autonomic testing identified varied profiles including POTS in 48\% of this selected sample.\par
\textit{Limit.} The figure is not a general-population prevalence and does not demonstrate RCCX, perinatal or \ensuremath{\alpha_2\delta} mediation.\par
\textit{Rival.} Referral selection and comorbidity explain the high clinical frequency.\par
\textit{Discriminator.} Population-based cohorts fail to reproduce an autonomic association under prespecified diagnostic criteria.\par
\textit{Sources.} \cite{ref33}\par
\subsection*{C26 · AI ideas \ensuremath{\rightarrow} individual story gains and collective similarity}
\noindent\textbf{Established · E4 · Randomized short-story writing experiment}\par
\textit{Claim.} Access to AI ideas improved average rated stories while AI-assisted stories were more similar to each other.\par
\textit{Limit.} The study does not establish benefit for scientific discovery, integrative phenotypes, or society-wide change.\par
\textit{Rival.} Task-specific anchoring explains similarity without a general AI effect.\par
\textit{Discriminator.} Preregistered replication on completed research artifacts finds no individual-quality or diversity contrast.\par
\textit{Sources.} \cite{ref21}\par
\subsection*{C27 · Severe CACNA2D1 loss \ensuremath{\rightarrow} impaired CaV2 function}
\noindent\textbf{Supported · E3 · Two rare human cases plus functional assays}\par
\textit{Claim.} Biallelic loss-of-function variants were linked to early epileptic encephalopathy and impaired CaV2 trafficking in assays.\par
\textit{Limit.} Rare severe variants do not identify effects of common variation or ordinary perinatal exposure.\par
\textit{Rival.} Other variants or pathways explain the clinical phenotype in these cases.\par
\textit{Discriminator.} Independent variant studies and orthogonal functional assays fail to reproduce the trafficking deficit.\par
\textit{Sources.} \cite{ref35}\par
